\subsection{UCI Classification}\label{app:uci-accuracies}

This section gives experimental details and additional results for the experiments evaluating generalization performance of the convex reformulations.

\textbf{Experimental Details}:
We selected 37 binary classification datasets from our filtered collection of 73 datasets;
see Appendix~\ref{app:uci-datasets} for details how the 73 datasets were obtained.

We used the default parameters for each of the convex solvers as described in Appendix~\ref{app:default-parameters},
except that a tighter convergence tolerance of \( 10^{-7} \) was used for terminating our methods.
R-FISTA was limited to \( 2000 \) iterations.
For the gated ReLU problems (both C-GReLU and NC-GReLU) we sampled the same set of \( 5000 \) activation patterns for both the convex reformulation and the original non-convex model.
We used \( 2500 \) activation patterns for the C-ReLU problem.

For each dataset-method pair, we performed five-fold cross validation on the training set
to select hyper-parameters.
We considered two hyper-parameters for our methods: regularization strength,
and the proportion of examples active in each local model (ie. the number of
non-zeros in each \( D_i \) matrix).
For regularization strength, we optimized over a logarithmic grid with values
	{ \small \(
		\cbr{1\times 10^{-8},
			3.59\times 10^{-8},
			1.29\times 10^{-7},
			4.64\times 10^{-7},
			1.67\times 10^{-6},
			5.99\times 10^{-6},
			2.15\times 10^{-5},
			7.74\times 10^{-5},
			2.78\times 10^{-4},
			1.0\times 10^{-3}}
		\)
	}.
For the proportion of active examples, we considered (1) setting the bias term
for each neuron to enforce \( 50 \% \) of examples to be active or (2)
setting the bias to \( 0 \) and allowing the proportion to be random.

For the baselines, we used the implementations available from
the \texttt{scikit-learn} package.
We optimized each random forest classifier with respect to the depth of
the random trees in the ensemble (\( \cbr{2, 4, 10, 25, 50}\))
and over the number of trees in the ensemble (\( \cbr{5, 10, 100, 1000} \)).
We used the standard soft-margin SVM with and chose regularization parameter from
the range
\( \cbr{1.\times 10^{-5}, 1.\times 10^{-4}, 1.\times 10^{-3}, 1.\times 10^{-2}, 1.\times 10^{-1}, 1} \)
for linear SVMs and
\( \cbr{1\times 10^{-5},
	1.78\times 10^{-4},
	3.16\times 10^{-3},
	5.62\times 10^{-2},
	1} \)
for SVMs with an RBF kernel.
For RBF SVMs, the RBF bandwidth was optimized over the grid \( \cbr{1\times 10^{-4},
	1.58\times 10^{-3},
	2.51\times 10^{-2},
	3.98\times 10^{-1},
	6.31,
	100.0} \).
To obtain final test accuracies, we re-trained each method on the full training
set.
For our methods, we report the best test accuracy out of five random restarts.

\begin{table}[t]
	\centering
	\caption{Additional results comparing our convex solvers against random forests (RF), SVMs with a linear kernel (Linear)
		and SVMs with an RBF kernel (RBF).
		We report test accuracies on a further 19 UCI datasets.
		Combined, C-GReLU and C-ReLU obtain the best test accuracy on 10 datasets.
		Out of the baselines considered, RBF SVMs are the most competitive with our approach,
		attaining or tying for best accuracy on 9 datasets.
	}%
	\label{table:binary-uci-accuracies-alt}
	\vspace{0.1in}
	\begin{small}
		\begin{tabular}{lccccc} \toprule
			\textbf{Dataset}  & \textbf{C-GReLU} & \textbf{C-ReLU} & \textbf{RF}   & \textbf{SVM}  & \textbf{RBF}  \\ \midrule
			breast-cancer     & 73.7             & 70.2            & \textbf{75.4} & 68.4          & 68.4          \\
			congressional     & 66.7             & 65.5            & 65.5          & \textbf{67.8} & \textbf{67.8} \\
			credit-approval   & 84.1             & 84.1            & \textbf{85.5} & \textbf{85.5} & 84.8          \\
			echocardiogram    & 80.8             & 76.9            & \textbf{88.5} & 84.6          & 84.6          \\
			haberman-survival & 67.2             & \textbf{75.4}   & 70.5          & 70.5          & 70.5          \\
			hepatitis         & 80.6             & 80.6            & \textbf{83.9} & 77.4          & 77.4          \\
			horse-colic       & 88.3             & 86.7            & \textbf{93.3} & 90.0          & \textbf{93.3} \\
			ionosphere        & 90.0             & 91.4            & 91.4          & 85.7          & \textbf{97.1} \\
			molec-biol        & 76.2             & \textbf{81.0}   & 76.2          & \textbf{81.0} & 66.7          \\
			monks-2           & \textbf{69.7}    & \textbf{69.7}   & 54.5          & 57.6          & \textbf{69.7} \\
			monks-3           & \textbf{95.8}    & \textbf{95.8}   & \textbf{95.8} & 87.5          & 91.7          \\
			musk-2            & 99.7             & \textbf{99.8}   & 97.3          & 95.1          & 99.6          \\
			parkinsons        & 97.4             & 97.4            & 84.6          & 89.7          & \textbf{100}  \\
			pittsburg         & \textbf{80.0}    & 75.0            & 75.0          & \textbf{80.0} & \textbf{80.0} \\
			ringnorm          & 97.4             & 83.0            & 95.1          & 76.9          & \textbf{98.2} \\
			spect             & 46.7             & 40.0            & 53.3          & \textbf{66.7} & \textbf{66.7} \\
			statlog-austr.    & 65.9             & \textbf{66.7}   & 62.3          & 65.2          & 65.2          \\
			statlog-heart     & 81.5             & \textbf{85.2}   & 83.3          & 81.5          & 81.5          \\
			twonorm           & 97.6             & \textbf{97.7}   & 97.2          & 97.4          & 97.4          \\
			vertebral-col.    & \textbf{91.9}    & \textbf{91.9}   & 87.1          & \textbf{91.9} & 90.3          \\ \bottomrule
		\end{tabular}
	\end{small}
\end{table}


\textbf{Additional Results}:
Table~\ref{table:binary-uci-accuracies-alt} reports test results for
19 of the 37 datasets, while Table~\ref{table:binary-uci-accuracies} in the
main paper presents results for the remaining 18 datasets.
Overall, we find that two-layer neural networks trained using our convex
solvers generally perform better than the baseline methods.
